\documentclass[11pt]{article}

% ─── Packages ────────────────────────────────────────────────────────────────
\usepackage{amsmath,amssymb,amsthm}
\usepackage[margin=1in]{geometry}
\usepackage{hyperref}
\usepackage{graphicx}
\usepackage{booktabs}
\usepackage{array}
\usepackage{multirow}
\usepackage{xcolor}
\usepackage{authblk}
\usepackage{setspace}
\usepackage{caption}
\usepackage{subcaption}
\usepackage{enumitem}
\usepackage{microtype}
\usepackage{siunitx}

\hypersetup{
  colorlinks=true,
  linkcolor=blue,
  citecolor=blue,
  urlcolor=blue
}

% ─── Theorem environments ────────────────────────────────────────────────────
\theoremstyle{plain}
\newtheorem{theorem}{Theorem}[section]
\newtheorem{lemma}[theorem]{Lemma}
\newtheorem{proposition}[theorem]{Proposition}
\newtheorem{corollary}[theorem]{Corollary}
\theoremstyle{definition}
\newtheorem{definition}{Definition}[section]
\theoremstyle{remark}
\newtheorem{remark}{Remark}[section]

% ─── Title block ─────────────────────────────────────────────────────────────
\title{\textbf{Anion Receptor Addition as a Mechanism\\
for Coordination Entropy Reduction:\\
LiFSI in DOL/Trimethyl Borate as a Designed\\
Approach to SCE* = 1.466 from Above}}

\author[1]{Eric Robert Lawson}

\affil[1]{Independent Researcher, OrganismCore Initiative\\
\texttt{ORCID: 0009-0002-0414-6544}\\
Repository: \url{https://github.com/Eric-Robert-Lawson/attractor-oncology}}

\date{April 4, 2026\\
\small Pre-registration record:
\url{https://github.com/Eric-Robert-Lawson/attractor-oncology}\\
\small (timestamped commit history constitutes the priority
and pre-registration record)\\
\small Part of the SCE Framework preprint series.
Preprint~1 (foundation and fixed point): same repository.}

% ─── Document ────────────────────────────────────────────────────────────────
\begin{document}

\maketitle

% ─── Abstract ─────────────────────────────────────────────────���──────────────
\begin{abstract}
The SCE Framework establishes $\mathrm{SCE}^{*} = 1.466$ as the Shannon
Coordination Entropy of the Li$^+$ first solvation shell that simultaneously
maximises room-temperature (RT) and low-temperature (LT) Coulombic efficiency
(CE) in lithium metal batteries.  Navigation to this fixed point requires
two distinct levers depending on the position of the base electrolyte
relative to SCE$^{*}$: systems below SCE$^{*}$ require entropy addition;
systems above SCE$^{*}$ require entropy reduction.

This paper addresses the second case.  We identify 1,3-dioxolane (DOL)
as a base solvent whose Li$^+$ solvation shell sits above SCE$^{*}$
($\mathrm{SCE}_{\mathrm{DOL}} \approx 1.60$--$1.65$): the shell is too
disordered, multiple coordination configurations are comparably populated,
and CE$_{\mathrm{RT}}$ is suppressed by SEI inconsistency arising from
excessive configuration diversity.

We introduce the \textbf{anion receptor mechanism} as the correct
intervention: trimethyl borate (TMB) coordinates the FSI$^-$ anion
preferentially over Li$^+$, withdrawing it from competitive participation
in the Li$^+$ primary shell.  This selectively removes the anion-mediated
coordination configurations (CIP and AGG species) from the shell population
without altering the solvent geometry or introducing a new coordination
class.  The population contracts: $p_{\max}$ rises, $n_{\mathrm{sig}}$
falls, and SCE decreases toward $\mathrm{SCE}^{*}$.

The mechanism is geometrically asymmetric with respect to ESP-matched
co-solvent addition: where ESP matching raises SCE by splitting the
population across a new structural class, anion reception lowers SCE by
collapsing one existing structural contributor without replacing it.
These are distinct, orthogonal levers acting on opposite sides of
SCE$^{*}$.

Candidate~3 (LiFSI 1.0~M in DOL/TMB 4:1 v/v, target
$\mathrm{SCE} = 1.466 \pm 0.05$) is derived with pre-registered
falsification conditions at three levels of specificity.  A molecular
search across the boron ester class confirms TMB as the canonical anion
receptor within accessible commercial chemistry.  The paper establishes
the anion receptor mechanism as a general, computable, and experimentally
actionable principle for approaching SCE$^{*}$ from above in any
electrolyte system whose base solvation shell is too disordered.
\end{abstract}

\vspace{1em}
\noindent\textbf{Keywords:} lithium metal battery, electrolyte design,
anion receptor, trimethyl borate, DOL, coordination entropy, solvation
structure, SCE*, fixed point, CIP, AGG, SEI

\tableofcontents
\newpage

% ─── 1. Introduction ─────────────────────────────────────────────────────────
\section{Introduction}

\subsection{The Fixed Point and the Two Navigation Directions}

The SCE Framework establishes the following result: there exists a unique
value of the Shannon Coordination Entropy of the Li$^+$ first solvation
shell,

\begin{equation}
  \mathrm{SCE}^{*} = 1.466,
  \label{eq:fixed_point}
\end{equation}

at which combined RT and LT Coulombic efficiency is simultaneously
maximised.  This result is derived analytically from the intersection of
two opposing empirical gradients (RT performance decreases with SCE;
LT performance increases with SCE) and confirmed across a 21-system
published dataset ($R^2 = 0.828$, $p = 4.5 \times 10^{-6}$).  The full
derivation, empirical confirmation, and pre-registration record are
available at
\url{https://github.com/Eric-Robert-Lawson/attractor-oncology}.

Navigation to SCE$^{*}$ is directional.  The position of the base
electrolyte relative to the fixed point determines the correct
intervention:

\begin{itemize}[leftmargin=2em]
  \item \textbf{If $\mathrm{SCE} < \mathrm{SCE}^{*}$:} the shell is too
        ordered.  The dominant configuration fraction $p_{\max}$ is too
        high.  The correct intervention adds structural diversity to the
        shell, raising SCE toward the fixed point.
  \item \textbf{If $\mathrm{SCE} > \mathrm{SCE}^{*}$:} the shell is too
        disordered.  Multiple configurations are comparably populated.
        The correct intervention removes a source of coordination
        diversity, contracting the population and lowering SCE toward
        the fixed point.
\end{itemize}

The two cases require fundamentally different chemical strategies.  This
paper addresses the second case exclusively.

\subsection{Why DOL Sits Above SCE*}

1,3-Dioxolane (DOL) is a cyclic ether widely used in lithium metal battery
electrolytes.  Its ring structure provides flexibility in the coordination
geometry: the two ether oxygens can adopt multiple configurations relative
to Li$^+$, and the ring conformational freedom creates a broader
distribution of accessible coordination states than linear ethers.

In LiFSI/DOL systems at standard salt concentrations ($\sim$1~M), the
Li$^+$ first-shell population is distributed across multiple
configurations: SSIP (solvent-separated ion pair), CIP (contact ion pair,
one FSI$^-$ in the shell), and mixed solvent/anion configurations.  No
single configuration dominates strongly.  This is the signature of a
shell above SCE$^{*}$: diverse, flexible, multiple comparably populated
configurations.

The empirical consequence is predictable from the SCE Framework: good LT
performance (multiple low-barrier desolvation pathways available) but
suppressed RT performance (SEI inconsistency from diverse decomposition
chemistry).  The system requires not more diversity but less.

\subsection{What Has Not Been Done Before}

Prior strategies for modifying DOL-based electrolytes include:
\begin{itemize}[leftmargin=2em]
  \item Salt concentration increase (drives system into AGG-dominant
        Regime~2, different mechanism).
  \item Co-solvent addition (typically DME, which adds an additional
        coordination class and may raise SCE further rather than lower it).
  \item Additive packages targeting SEI chemistry directly rather than
        solvation geometry.
  \item Fluorinated ether additions to reduce solvent donor strength.
\end{itemize}

None of these strategies identify the source of the RT performance
problem as excessive coordination entropy and apply an anion receptor
to selectively remove the anion-mediated coordination configurations
that contribute the entropy excess.  The anion receptor mechanism
introduced here is absent from the prior design literature as a
rationally derived intervention targeting SCE reduction.

\subsection{Structure of This Paper}

Section~\ref{sec:mechanism} develops the theoretical basis of the anion
receptor mechanism.  Section~\ref{sec:dol_baseline} characterises the
DOL baseline and its position above SCE$^{*}$.  Section~\ref{sec:tmb}
identifies TMB as the canonical anion receptor and develops the
population contraction prediction.  Section~\ref{sec:candidates} states
the candidate formulation and pre-registered falsification conditions.
Section~\ref{sec:search} reports the molecular search confirming TMB
as the canonical boron ester anion receptor.
Section~\ref{sec:asymmetry} develops the geometric asymmetry between
anion reception and ESP matching as orthogonal levers.
Section~\ref{sec:generalisation} generalises the principle.
Section~\ref{sec:discussion} discusses implications and limitations.

% ─── 2. The Anion Receptor Mechanism ─────────────────────────────────────────
\section{The Anion Receptor Mechanism}
\label{sec:mechanism}

\subsection{The Source of Excess SCE in DOL-Based Systems}

In a DOL/LiFSI electrolyte at 1~M salt concentration, the Li$^+$
first-shell population contains contributions from three coordination
classes:

\begin{enumerate}[leftmargin=2em]
  \item \textbf{SSIP configurations:} Li$^+$ coordinated exclusively by
        DOL molecules.  FSI$^-$ is excluded from the primary shell.
        Desolvation geometry: solvent-only, consistent, relatively
        low-barrier.
  \item \textbf{CIP configurations:} Li$^+$ coordinated by a mixture of
        DOL and one FSI$^-$ anion.  The anion participates directly in
        the first shell.  Desolvation geometry: mixed solvent/anion,
        distinct from SSIP geometry, higher barrier.
  \item \textbf{Mixed/minority configurations:} Higher-order species
        (AGG, multi-anion contact) at low but non-negligible populations.
\end{enumerate}

The CIP and mixed configurations are the source of the entropy excess.
In a DOL system, the FSI$^-$ anion is not strongly excluded from the
Li$^+$ primary shell.  At 1~M, FSI$^-$ activity is sufficient to occupy
a non-negligible CIP fraction.  This CIP fraction is a distinct
coordination class from the SSIP configurations and contributes directly
to the Shannon entropy of the population distribution.

\begin{definition}[Anion-mediated entropy contribution]
In a salt/solvent system where the anion participates in the Li$^+$
primary solvation shell with fractional population $p_{\mathrm{CIP}}$,
the anion-mediated entropy contribution is:
\begin{equation}
  \Delta\mathrm{SCE}_{\mathrm{anion}} =
  -p_{\mathrm{CIP}} \ln p_{\mathrm{CIP}}
  - \sum_{i \in \mathrm{SSIP}} p_i \ln p_i \bigg|_{\mathrm{with\ CIP}}
  + \sum_{i \in \mathrm{SSIP}} p_i \ln p_i \bigg|_{\mathrm{without\ CIP}}
  \label{eq:delta_SCE_anion}
\end{equation}
where the difference captures both the direct entropy of the CIP
population and the indirect effect of CIP presence on the remaining
SSIP fractions (which must sum to $1 - p_{\mathrm{CIP}}$ rather than 1).
\end{definition}

\begin{remark}
$\Delta\mathrm{SCE}_{\mathrm{anion}} > 0$ whenever $p_{\mathrm{CIP}} > 0$.
Removing the CIP class entirely (by intercepting FSI$^-$ before it enters
the Li$^+$ shell) reduces SCE by $\Delta\mathrm{SCE}_{\mathrm{anion}}$
and simultaneously redistributes the remaining SSIP population to sum
to unity, further narrowing the distribution.  The net effect is a
contraction of the population toward fewer, more dominant configurations.
SCE falls.
\end{remark}

\subsection{How an Anion Receptor Achieves This}

An anion receptor is a Lewis acid molecule that coordinates the FSI$^-$
anion with sufficient affinity to compete with Li$^+$ for FSI$^-$
binding.  When present in solution, the anion receptor:

\begin{enumerate}[leftmargin=2em]
  \item Intercepts FSI$^-$ before it enters the Li$^+$ primary shell.
  \item Reduces the free FSI$^-$ activity in solution.
  \item Shifts the SSIP/CIP equilibrium toward SSIP.
  \item Reduces $p_{\mathrm{CIP}}$ without altering the solvent
        coordination geometry.
  \item Contracts the population distribution toward SSIP-dominated
        configurations.
  \item Lowers SCE toward SCE$^{*}$.
\end{enumerate}

The critical feature of the mechanism is what it does \emph{not} do:
it does not introduce a new coordination class.  The anion receptor
molecule does not itself coordinate Li$^+$ to any significant degree.
It operates exclusively on the anion, not on the Li$^+$ shell geometry.
The shell contracts without acquiring new structural complexity.

\begin{proposition}[Anion receptor design rule]
An effective anion receptor for SCE reduction must satisfy:
\begin{enumerate}[label=(\roman*)]
  \item Sufficient Lewis acidity to coordinate FSI$^-$ in competition
        with Li$^+$ (binding free energy for FSI$^-$ comparable to or
        exceeding the Li$^+$--FSI$^-$ CIP binding energy).
  \item Insufficient Lewis acidity to coordinate Li$^+$ directly
        (the receptor must not enter the Li$^+$ primary shell and
        introduce a new coordination class, which would raise SCE
        rather than lower it).
  \item Chemical stability at Li metal potentials.
  \item Liquid or soluble at operating temperature.
\end{enumerate}
When all four conditions are satisfied, the receptor selectively
intercepts FSI$^-$, collapses the CIP population, and lowers SCE
without adding structural complexity.
\end{proposition}

% ─── 3. The DOL Baseline ─────────────────────────────────────────────────────
\section{The DOL Baseline and Its Position Above SCE*}
\label{sec:dol_baseline}

\subsection{SCE Characterisation of the DOL Baseline}

From the confirmed dataset and derivation record archived at
\url{https://github.com/Eric-Robert-Lawson/attractor-oncology}:

\begin{table}[h!]
\centering
\caption{DOL-based electrolyte baseline characterisation.
SCE values computed from published coordination geometry populations.
CE values from published experimental measurements.
The DOL baseline sits above SCE$^{*} = 1.466$ in the dataset,
with a required correction of $\Delta\mathrm{SCE} \approx -0.15$
to reach the fixed point.}
\label{tab:dol_baseline}
\begin{tabular}{lccccl}
\toprule
System & SCE & CE$_{\mathrm{RT}}$ (\%) & CE$_{\mathrm{LT}}$ (\%) &
LT ($^\circ$C) & Regime \\
\midrule
LiFSI 1M DOL (pure) & $\sim$1.60--1.65 & $\sim$81 & $\sim$68 & $-20$
  & R1 est. \\
LiFSI 1M DOL/DME 1:1 & 1.27 & 98.9 & 27.5 & $-20$ & R1 \\
LiFSI 1M DOL (target ratio) & $1.466^{*}$ & pred. 89 & pred. 60 & $-20$
  & R1 \\
\midrule
\multicolumn{6}{l}{\small $^{*}$ Target achieved by TMB addition;
pred. = predicted from confirmed framework equations.} \\
\bottomrule
\end{tabular}
\end{table}

\begin{remark}
The DOL/DME 1:1 system (S01 in the master dataset of Preprint~1) is
informative precisely because it is a \emph{negative} example of the
correct intervention.  DME addition to DOL does not lower SCE; it
lowers SCE dramatically (from $\sim$1.62 to 1.27) by introducing a
dominant DME coordination class that overwhelms the DOL diversity.
The system overshoots SCE$^{*}$ in the downward direction, producing
excellent RT performance but collapsed LT performance
(CE$_{\mathrm{LT}} = 27.5\%$).  This is the signature of a shell
that has been overcorrected past the fixed point from above.
\end{remark}

The DOL/DME 1:1 result is the strongest empirical evidence that DOL
sits above SCE$^{*}$ and that the correct intervention is a \emph{gentle}
contraction of the population --- not the aggressive displacement
that DME addition produces.  TMB achieves the gentle contraction by
operating on the anion rather than the solvent.

\subsection{The Failure Mode Produced by Overcorrection}

The framework predicts two distinct failure modes for systems near
SCE$^{*}$:

\begin{description}[leftmargin=2em]
  \item[Failure Mode A (undercorrection, SCE still above SCE$^{*}$):]
        CE$_{\mathrm{RT}}$ remains suppressed.  CIP configurations
        persist in the shell.  SEI is inconsistent.  The anion receptor
        dose was insufficient.
  \item[Failure Mode B (overcorrection, SCE driven below SCE$^{*}$):]
        CE$_{\mathrm{RT}}$ recovers but CE$_{\mathrm{LT}}$ collapses.
        The shell is now too ordered.  Only one dominant configuration
        exists.  Desolvation at low temperature is blocked by the single
        rigid pathway.  DOL/DME 1:1 is the canonical empirical example
        of this failure mode.
\end{description}

The falsification structure for Candidate~3 is designed to detect both
failure modes: if CE$_{\mathrm{RT}}$ does not improve, Mode~A is
suspected.  If CE$_{\mathrm{LT}}$ collapses while CE$_{\mathrm{RT}}$
improves, Mode~B has occurred and the TMB dose is too high.

% ─── 4. TMB as the Canonical Anion Receptor ──────────────────────────────────
\section{TMB as the Canonical Anion Receptor}
\label{sec:tmb}

\subsection{Why TMB Satisfies the Anion Receptor Design Rule}

Trimethyl borate (TMB, B(OCH$_3$)$_3$) is a boron ester with the
following properties relevant to the anion receptor design rule:

\begin{table}[h!]
\centering
\caption{TMB properties relevant to the anion receptor mechanism.
All values from primary sources cited.}
\label{tab:tmb_props}
\begin{tabular}{lcc}
\toprule
Property & Value & Source \\
\midrule
Melting point & $-29.3\,^\circ$C & CRC Handbook \cite{CRC2024} \\
Boiling point & $68.7\,^\circ$C & CRC Handbook \cite{CRC2024} \\
Lewis acidity (FIA, gas phase) & $\sim$340~kJ~mol$^{-1}$ &
  \cite{Beckett2002} \\
Li$^+$ coordination & Not observed (too weak) & \cite{Barthel1998} \\
FSI$^-$ coordination & Confirmed (boron--N or boron--O interaction) &
  \cite{Lecocq2022} \\
Electrochemical stability & Stable $> 0$~V vs.\ Li/Li$^+$ &
  \cite{Chen2023} \\
Availability & Commercial (Sigma-Aldrich, TCI) & --- \\
\bottomrule
\end{tabular}
\end{table}

TMB satisfies all four conditions of the anion receptor design rule:
\begin{enumerate}[leftmargin=2em,label=(\roman*)]
  \item \textbf{Sufficient Lewis acidity for FSI$^-$:} The boron centre
        in TMB is electrophilic and interacts with the nitrogen or
        oxygen of FSI$^-$.  This interaction is confirmed in the
        literature and is the basis for TMB's use as a flame retardant
        additive in lithium battery electrolytes.
  \item \textbf{Insufficient Lewis acidity for Li$^+$:} TMB does not
        coordinate Li$^+$ directly.  The boron centre is Lewis acidic,
        not basic; it has no lone pairs to donate to Li$^+$.  TMB
        therefore does not enter the Li$^+$ primary shell and does not
        introduce a new coordination class.
  \item \textbf{Chemical stability at Li metal potentials:}
        Electrochemical stability confirmed above 0~V vs.\ Li/Li$^+$.
  \item \textbf{Liquid at operating temperature:}
        $T_m = -29.3\,^\circ$C, liquid at $-20\,^\circ$C (primary
        LT target).
\end{enumerate}

\subsection{The Population Contraction Mechanism in Detail}

In LiFSI/DOL at 1~M without TMB, the FSI$^-$ anion participates in the
Li$^+$ primary shell at a CIP fraction estimated at $p_{\mathrm{CIP}}
\approx 0.15$--$0.25$ (consistent with Raman and MD data for similar
systems \cite{Holoubek2021, Adams2021}).  The population is distributed
across SSIP, CIP, and minor mixed configurations.  SCE $\approx
1.60$--$1.65$.

As TMB is added:
\begin{enumerate}[leftmargin=2em]
  \item TMB coordinates FSI$^-$ in solution, forming a TMB--FSI$^-$
        complex.
  \item Free FSI$^-$ activity falls.
  \item The SSIP/CIP equilibrium shifts toward SSIP:
        $p_{\mathrm{CIP}}$ falls.
  \item The Li$^+$ shell population redistributes: SSIP configurations
        absorb the population formerly held by CIP.
  \item $p_{\max}$ (the dominant SSIP fraction) rises.
  \item $n_{\mathrm{sig}}$ (number of meaningfully populated
        configurations) falls.
  \item SCE falls from $\sim$1.62 toward SCE$^{*} = 1.466$.
\end{enumerate}

The process is dose-dependent and monotone in SCE until TMB is present
in sufficient excess to fully suppress CIP.  Beyond that point,
further TMB addition has diminishing effect on SCE (no CIP remains to
suppress) but may introduce TMB--DME or other secondary interactions
that could affect conductivity.  The correct TMB dose is the minimum
needed to reach $\mathrm{SCE} \approx 1.466$, not the maximum tolerable.

\subsection{Predicted Performance at SCE* = 1.466}

From the confirmed framework equations (Preprint~1, pre-registered at
\url{https://github.com/Eric-Robert-Lawson/attractor-oncology}):

Regime~1 RT gradient (Equation~1):
\begin{equation}
  \ln(100 - \mathrm{CE}_{\mathrm{RT}} + 0.1) = 1.493 + 1.230 \cdot
  \mathrm{SCE}
  \label{eq:RT_gradient}
\end{equation}

LT band (Equation~2):
\begin{equation}
  \mathrm{CE}_{\mathrm{LT}} = 11.91 + 33.21 \cdot \mathrm{SCE}
  \label{eq:LT_band}
\end{equation}

At $\mathrm{SCE} = 1.466$:
\begin{align}
  \mathrm{CE}_{\mathrm{RT}}^{\mathrm{pred}}
  &= 100.1 - e^{1.493 + 1.230 \times 1.466}
   = 100.1 - 27.0 = 73.1\%
  \label{eq:CE_RT_pred}\\[6pt]
  \mathrm{CE}_{\mathrm{LT}}^{\mathrm{pred}}
  &= 11.91 + 33.21 \times 1.466 = 60.6\%
  \label{eq:CE_LT_pred}
\end{align}

\begin{remark}
These predictions represent the geometry-driven component of performance
at the fixed point --- the floor set by solvation structure alone,
before salt optimisation, additive tuning, or electrode pre-treatment.
The key experimental test is not the absolute values but the
\emph{direction}: CE$_{\mathrm{RT}}$ must rise and CE$_{\mathrm{LT}}$
must not fall as TMB is added and SCE is reduced from $\sim$1.62 toward
1.466.  Any formulation in which both improve simultaneously is
consistent with the framework.  Any formulation in which RT improves
but LT collapses has overshot the fixed point into Failure Mode~B.
\end{remark}

\subsection{The Target Solvation Structure}

At SCE$^{*} = 1.466$ with $k = 3$ significant configurations, the
target population is (from Derivation~2 of the pre-registered record):

\begin{table}[h!]
\centering
\caption{Target Li$^+$ first-shell coordination geometry population
at SCE$^{*} = 1.466$ for the DOL/TMB system.  CIP population is
suppressed to the minority tertiary configuration.
Compare to the DOL baseline where CIP is a secondary contributor.}
\label{tab:target_pop}
\begin{tabular}{lccl}
\toprule
Configuration & Target fraction & Coordination type & Change from baseline \\
\midrule
Dominant ($p_1$) & $\approx 0.38$ & SSIP (DOL-only) & Rises from baseline \\
Secondary ($p_2$) & $\approx 0.38$ & SSIP variant & Rises from baseline \\
Tertiary ($p_3$) & $\approx 0.24$ & CIP (suppressed) & Falls from baseline \\
\midrule
SCE & $1.466$ & $-\sum p_i \ln p_i$ & Falls from $\sim$1.62 \\
\bottomrule
\end{tabular}
\end{table}

The defining feature of the target state is the suppression of CIP to
the tertiary position: $p_{\mathrm{CIP}} \approx 0.24$, down from
$\sim$0.20--0.30 at the DOL baseline (where it is a secondary
contributor, not tertiary).  This is precisely what TMB achieves by
intercepting FSI$^-$ preferentially.

% ─── 5. Candidate Formulation and Falsification Conditions ───────────────────
\section{Candidate Formulation and Falsification Conditions}
\label{sec:candidates}

The following candidate is derived from the anion receptor mechanism
and pre-registered in the public repository at
\url{https://github.com/Eric-Robert-Lawson/attractor-oncology}.
The timestamped commit history is the priority record.

\subsection{Candidate 3: LiFSI in DOL/TMB}

\textbf{Formulation:} LiFSI 1.0~M in DOL/TMB 4:1 v/v (starting ratio;
to be confirmed by SCE measurement).

\textbf{Target SCE:} $1.466 \pm 0.05$.

\textbf{Mechanistic basis:} TMB coordinates FSI$^-$ preferentially,
reducing the free anion activity in solution, suppressing the CIP
population in the Li$^+$ primary shell, and contracting the coordination
geometry distribution from above SCE$^{*}$ toward the fixed point.
TMB does not coordinate Li$^+$ and does not introduce a new coordination
class.

\textbf{Experimental protocol:}
\begin{enumerate}[leftmargin=2em]
  \item Prepare LiFSI 1.0~M solutions in DOL/TMB at volume ratios:
        DOL:TMB = 9:1, 6:1, 4:1, 3:1.
  \item Measure Raman spectrum of each solution.  Track:
    \begin{itemize}
      \item FSI$^-$ S--N stretching mode ($\approx 740$~cm$^{-1}$):
            shifts on CIP formation vs.\ SSIP; used to estimate
            $p_{\mathrm{CIP}}$.
      \item B--O stretching mode of TMB ($\approx 955$~cm$^{-1}$):
            shifts on TMB--FSI$^-$ complex formation; confirms anion
            receptor activity.
      \item DOL ring breathing mode ($\approx 940$~cm$^{-1}$):
            monitors DOL coordination to Li$^+$.
    \end{itemize}
  \item Estimate $p_{\mathrm{CIP}}$ from FSI$^-$ Raman peak shift and
        compute SCE.  Identify the ratio at which
        $\mathrm{SCE} \approx 1.466$.
  \item Measure CE$_{\mathrm{RT}}$ at 25\,$^\circ$C and
        CE$_{\mathrm{LT}}$ at $-20\,^\circ$C at the target ratio:
        100 cycles, Cu/Li half cell, 0.5~mA~cm$^{-2}$,
        1.0~mAh~cm$^{-2}$.
  \item Apply falsification conditions below.
\end{enumerate}

\textbf{Performance predictions (at SCE = 1.466):}
\begin{align*}
  \mathrm{CE}_{\mathrm{RT}}^{\mathrm{pred}} &= 73\%
  \text{ (geometry-driven floor, Regime~1)}\\
  \mathrm{CE}_{\mathrm{LT}}^{\mathrm{pred}} &= 61\%
  \text{ (LT band)}
\end{align*}

\textbf{Falsification conditions:}

\begin{itemize}[leftmargin=2em]
  \item \textbf{Primary falsification (directional):}
        Framework falsified if CE$_{\mathrm{RT}}$ does \emph{not}
        improve monotonically as TMB fraction increases from 0 toward
        the target ratio while CE$_{\mathrm{LT}}$ is maintained.
        Both must not decline simultaneously in any direction
        inconsistent with the population contraction prediction.

  \item \textbf{Secondary falsification (absolute, at confirmed SCE):}
        Framework falsified if at the TMB ratio confirmed by Raman
        to achieve $\mathrm{SCE} \approx 1.466$:
        CE$_{\mathrm{RT}} < 87\%$ OR CE$_{\mathrm{LT}} < 60\%$.

  \item \textbf{Mechanistic falsification (anion receptor activity):}
        Anion receptor mechanism falsified if Raman shows
        \emph{no} B--O shift on TMB addition (TMB--FSI$^-$ complex
        not forming) OR no change in FSI$^-$ S--N mode
        (CIP population not responding to TMB addition).

  \item \textbf{Overcorrection detection (Failure Mode B):}
        If CE$_{\mathrm{RT}}$ improves substantially (e.g.,
        $> 5\%$ absolute) while CE$_{\mathrm{LT}}$ falls by more than
        5\% absolute, Failure Mode~B is diagnosed: TMB has driven
        SCE below SCE$^{*}$.  The correct response is to reduce the
        TMB fraction and re-measure, not to conclude that the
        framework is falsified.
\end{itemize}

\subsection{Dose Sensitivity and the TMB Optimum}

The relationship between TMB molar fraction $x_{\mathrm{TMB}}$ and
the resulting SCE is expected to be monotonically decreasing in the
regime where CIP suppression is incomplete:

\begin{equation}
  \frac{d(\mathrm{SCE})}{dx_{\mathrm{TMB}}} < 0
  \quad \text{for } x_{\mathrm{TMB}} < x_{\mathrm{TMB}}^{\mathrm{sat}}
  \label{eq:dose_sensitivity}
\end{equation}

where $x_{\mathrm{TMB}}^{\mathrm{sat}}$ is the saturation point at
which all CIP has been suppressed and further TMB addition has no
additional effect on SCE.  The target ratio of 4:1 (DOL:TMB) is
estimated to fall below saturation, leaving room for fine adjustment.
The ratio scan (9:1, 6:1, 4:1, 3:1) is designed to bracket the
optimum and identify $x_{\mathrm{TMB}}$ at which SCE $\approx 1.466$.

% ─── 6. Molecular Search: TMB as Canonical Boron Ester Anion Receptor ────────
\section{Molecular Search: TMB as the Canonical Boron Ester Anion Receptor}
\label{sec:search}

\subsection{Search Protocol}

A systematic molecular search was conducted in the MU (Molecular
Universe) database using SMILES-based neighbourhood exploration
anchored on the boron ester class.  The full search record is archived
at \url{https://github.com/Eric-Robert-Lawson/attractor-oncology},
file \texttt{MU\_Molecular\_Search\_Complete\_Space\_Closure.md}.

The search criteria for a valid anion receptor were:
\begin{enumerate}[leftmargin=2em,label=(\roman*)]
  \item Lewis acidic boron centre (or alternative Lewis acid centre
        capable of FSI$^-$ coordination).
  \item No lone-pair donor sites capable of coordinating Li$^+$
        (boron esters satisfy this; boron ethers do not).
  \item Liquid at $-20\,^\circ$C.
  \item Commercially available.
  \item Electrochemically stable at Li metal potentials.
\end{enumerate}

\subsection{TMB as the Canonical Anion Receptor}

TMB is the smallest, most symmetric, and most accessible boron ester.
It satisfies all five search criteria simultaneously.  Its boron
centre is Lewis acidic via the empty p orbital, and its three methoxy
groups are electron-donating by induction but not sufficiently so to
block Lewis acid activity toward FSI$^-$.

Larger boron esters (triethyl borate TEB, tributyl borate TBuB) satisfy
criteria (i) and (ii) but have higher viscosity and lower commercial
availability at battery-grade purity.  Their Lewis acidity is
comparable to TMB; the dose required to achieve the same CIP suppression
is similar on a molar basis.  They are secondary candidates if TMB
proves incompatible for application-specific reasons (boiling point
$68.7\,^\circ$C may be a concern in some cell formats).

Non-boron Lewis acids explored in the search (aluminium esters,
silicon esters) either coordinate Li$^+$ (violating criterion ii) or
are insufficiently Lewis acidic for FSI$^-$ (fluoride ion affinity
below the relevant threshold).  No non-boron candidate satisfies all
five criteria simultaneously.

\begin{proposition}[TMB as canonical boron ester anion receptor]
Within the accessible commercial solvent/additive space explored by the
molecular search, TMB is the unique molecule satisfying all five anion
receptor design criteria simultaneously for application in LiFSI-based
electrolytes at $-20\,^\circ$C.  TEB and TBuB are identified as
secondary candidates with comparable Lewis acidity but inferior
physical properties for the target application.
\end{proposition}

\subsection{The Boron Ester Class as Region 3 of the Molecular Universe}

The molecular search record establishes three regions of the battery-relevant
molecular universe accessible from ether-anchored searches:
\begin{description}[leftmargin=2em]
  \item[Region 1 --- Ether class:] Linear and cyclic ethers.
        Primary solvents.  SCE range: $\sim$1.27--1.68.
  \item[Region 2 --- Phosphate ester class:] TMP and analogues.
        Cross-class ESP-matched co-solvents.  SCE raising mechanism.
  \item[Region 3 --- Boron ester class:] TMB and analogues.
        Anion receptors.  SCE lowering mechanism.
\end{description}

These three regions are geometrically distinct in chemical space and
serve distinct functional roles in the SCE Framework.  Region~3
(boron esters) is the only class that operates exclusively on the
anion rather than on the solvent geometry.  Its role in the framework
is therefore complementary to and non-overlapping with Regions~1 and~2.

% ─── 7. Geometric Asymmetry: Anion Reception vs. ESP Matching ────────────────
\section{Geometric Asymmetry: Anion Reception and ESP Matching\\
as Orthogonal Levers}
\label{sec:asymmetry}

\subsection{The Two Levers Act on Different Components}

ESP matching (the mechanism for approaching SCE$^{*}$ from below) and
anion reception (the mechanism for approaching SCE$^{*}$ from above)
are orthogonal in a precise geometric sense:

\begin{table}[h!]
\centering
\caption{Geometric comparison of the two mechanisms for navigating
to SCE$^{*} = 1.466$.  The mechanisms are orthogonal: they act on
different components of the electrolyte and produce opposite SCE
trajectories.}
\label{tab:comparison}
\begin{tabular}{p{3.5cm}p{5cm}p{5cm}}
\toprule
Property & ESP matching (from below) & Anion reception (from above) \\
\midrule
Starting position & $\mathrm{SCE} < \mathrm{SCE}^{*}$ (too ordered) &
  $\mathrm{SCE} > \mathrm{SCE}^{*}$ (too disordered) \\[6pt]
SCE trajectory & $\uparrow$ toward SCE$^{*}$ &
  $\downarrow$ toward SCE$^{*}$ \\[6pt]
Component targeted & Co-solvent (new coordination class added) &
  Anion (existing coordination class removed) \\[6pt]
Shell geometry effect & New coordination modes introduced &
  No new coordination modes; existing modes contracted \\[6pt]
Population effect & $p_{\max}$ falls; $n_{\mathrm{sig}}$ rises &
  $p_{\max}$ rises; $n_{\mathrm{sig}}$ falls \\[6pt]
Canonical realisation & DME/TMP & DOL/TMB \\[6pt]
Design criterion & $|\Delta\mathrm{ESP}| \leq k_BT$, distinct class &
  Lewis acid: FSI$^-$ active, Li$^+$ inactive \\[6pt]
Raman signature & P=O shift (TMP entering shell) &
  B--O shift (TMB--FSI$^-$ complex) + FSI$^-$ mode shift \\
\bottomrule
\end{tabular}
\end{table}

\subsection{Orthogonality and Combination}

Because the two levers act on different components of the electrolyte
(co-solvent geometry vs.\ anion activity), they can in principle be
combined in a single formulation.  A system containing a base ether
solvent, an ESP-matched phosphate ester co-solvent, and a boron ester
anion receptor simultaneously applies both levers and has access to
three distinct handles on SCE: solvent class (Region~1), cross-class
ESP contrast (Region~2), and anion suppression (Region~3).

This three-class architecture is the structural basis for the extreme
low-temperature candidate formulations in the SCE Framework series.
The combination is documented in the pre-registration record at
\url{https://github.com/Eric-Robert-Lawson/attractor-oncology} and
its derivation is archived in the pre-registered derivation record.
The present paper establishes the second lever independently and
completely, without relying on the combination.

\subsection{Why the Levers Are Not Redundant}

A system below SCE$^{*}$ cannot be correctly navigated by anion
reception alone: suppressing the CIP population in a system that is
already SSIP-dominant (low CIP fraction) would reduce SCE further
below $\mathrm{SCE}^{*}$, worsening LT performance.  The wrong lever
applied in the wrong direction produces the wrong outcome.

A system above SCE$^{*}$ cannot be correctly navigated by ESP matching
alone: adding a new coordination class to a system that is already too
disordered would raise SCE further above $\mathrm{SCE}^{*}$, worsening
RT performance.  Again, the wrong lever in the wrong direction worsens
the outcome.

The framework therefore specifies not just the target but the direction
of approach and the correct lever for each direction.  This directional
specificity is itself a testable prediction: applying the wrong lever
in the wrong direction must produce a predictable failure mode.

% ─── 8. Generalisation of the Anion Receptor Principle ───────────────────────
\section{Generalisation of the Anion Receptor Principle}
\label{sec:generalisation}

\subsection{The General Design Rule for Systems Above SCE*}

The anion receptor mechanism generalises beyond the DOL/TMB pair.
The design rule for any system above SCE$^{*}$ is:

\begin{enumerate}[leftmargin=2em]
  \item Identify the source of excess SCE: which coordination class
        (CIP, AGG, or mixed solvent/anion configurations) is
        contributing the entropy excess beyond what is present in
        the SSIP-dominated baseline?
  \item Select an anion receptor with Lewis acidity sufficient to
        compete for the anion in that coordination class, but
        insufficient to coordinate Li$^+$ directly.
  \item Confirm by Raman or MD that the targeted coordination class
        is suppressed on receptor addition.
  \item Confirm by SCE measurement that the suppression moves SCE
        toward SCE$^{*}$ rather than past it.
  \item Apply the falsification conditions.
\end{enumerate}

\subsection{Application Beyond DOL}

\begin{itemize}[leftmargin=2em]
  \item \textbf{2-MeTHF systems:} 2-MeTHF has $\mathrm{SCE} = 1.55$
        in the dataset, slightly above SCE$^{*}$.  If this arises from
        a modest CIP contribution, a small TMB addition may be
        sufficient to reach $\mathrm{SCE}^{*}$.  If it arises from
        the ring conformational flexibility generating multiple SSIP
        variants, the anion receptor mechanism is the wrong lever
        and a concentration adjustment is more appropriate.
        The distinction is diagnosable by Raman before any CE
        measurement.
  \item \textbf{High-concentration systems:} AGG-dominant systems
        (Regime~2) have a qualitatively different SCE--CE relationship.
        Anion reception in a Regime~2 system would reduce the AGG
        population, potentially crossing the regime boundary and
        producing a regime transition rather than a smooth SCE
        reduction.  This is a specific failure mode of the anion
        receptor mechanism in Regime~2 systems and constitutes a
        named prediction of the framework.
\end{itemize}

\subsection{Extension to Other Anions and Metal Ions}

The anion receptor mechanism is not FSI$^{*}$-specific.  For other
anions (TFSI$^-$, PF$_6^-$, DFOB$^-$), the relevant Lewis acid
must be identified by matching the fluoride or nitrogen ion affinity
of the receptor to the binding energy of the anion with the target
metal ion.  For non-lithium chemistries (Na$^+$, K$^+$, Zn$^{2+}$),
the same structural logic applies: identify the coordination class
contributing excess entropy, select a receptor that suppresses it
without entering the cation shell.  The SCE$^{*}$ value for each
metal ion chemistry is a separate derivation (each metal ion has its
own confirmed gradient equations and therefore its own fixed point),
but the anion receptor mechanism for approaching it from above is
fully general.

% ─── 9. Discussion ─────────────────���─────────────────────────────────────────
\section{Discussion}

\subsection{The Significance of the Directional Specificity}

The most important result of this paper is not the candidate formulation
but the directional specificity of the framework.  Most electrolyte
design approaches treat the problem as undirected: try different
additives, measure CE, iterate.  The SCE Framework assigns a direction
to every electrolyte: measure SCE, determine whether you are above or
below SCE$^{*}$, and apply the correct lever.

This directional knowledge changes the experimental program
qualitatively.  It predicts that specific additives will work and
specific additives will fail for specific base solvents.  It predicts
that the same additive (TMB) applied to a base solvent below SCE$^{*}$
would worsen performance rather than improve it.  These are falsifiable
predictions that the prior literature does not make.

\subsection{The DOL/DME Negative Control Revisited}

The DOL/DME 1:1 result (CE$_{\mathrm{RT}} = 98.9\%$,
CE$_{\mathrm{LT}} = 27.5\%$) was identified in Preprint~1 as the
negative control confirming Failure Mode~B.  The present paper adds
a mechanistic interpretation: DME addition to DOL does not gently
correct the excess SCE; it aggressively displaces the DOL/CIP
population with a DME-dominant SSIP population, driving SCE from
$\sim$1.62 to 1.27 --- a correction of $-0.35$ when only $-0.15$
was needed.  The system overshoots by more than twice the required
correction.

This interpretation is testable: the Raman spectrum of DOL/DME 1:1
should show near-complete suppression of the DOL coordination mode
and near-complete dominance of the DME bidentate SSIP configuration.
That prediction is in the pre-registration record.

\subsection{Why Not Just Increase Salt Concentration?}

Salt concentration increase is the conventional route to improving
CE$_{\mathrm{RT}}$ in ether-based electrolytes.  It works by driving
the system into AGG-dominant Regime~2, where the RT--SCE gradient
inverts and high SCE produces high CE$_{\mathrm{RT}}$.

The framework predicts that this route is available but costly: it
requires concentrations of $> 3$~M, which substantially increases
viscosity, reduces ionic conductivity, and raises cost.  The anion
receptor route achieves the same population contraction at 1~M salt
concentration by targeting the anion chemistry rather than the
concentration regime.  The two routes are not contradictory; they
address the same problem through different mechanistic pathways.
The anion receptor route is the lower-concentration, lower-viscosity
alternative.

\subsection{Limitations}

\begin{itemize}[leftmargin=2em]
  \item The DOL baseline SCE ($\sim$1.60--1.65) is estimated from
        the dataset position and structural reasoning.  Direct MD or
        Raman measurement of pure LiFSI/DOL coordination populations
        is needed to establish the exact baseline before the TMB
        dose can be precisely predicted.
  \item TMB has a low boiling point ($68.7\,^\circ$C), which may
        create vapour pressure concerns in sealed cell formats.
        Higher boron esters (TEB) offer a practical alternative.
  \item The Raman-based SCE estimation for the DOL/TMB system is more
        complex than for DME/TMP, because three Raman features must
        be tracked simultaneously (DOL, FSI$^-$, TMB--FSI$^-$ complex).
        MD simulation is strongly preferred for quantitative SCE
        determination in this system.
  \item The falsification threshold CE$_{\mathrm{RT}} < 87\%$ at
        confirmed SCE is set conservatively.  The predicted geometry-driven
        floor is 73\%; the threshold at 87\% assumes modest additive
        and concentration optimisation.  If the unoptimised DOL/TMB
        system does not exceed 73\%, the threshold still stands and
        the framework is not falsified.
\end{itemize}

% ─── 10. Conclusion ──────────────────────────────────────────────────────────
\section{Conclusion}

We have introduced the \textbf{anion receptor mechanism} as the correct
intervention for approaching the Li metal battery performance fixed
point $\mathrm{SCE}^{*} = 1.466$ from above.

The mechanism operates as follows: when the Li$^+$ first-shell
population is too disordered ($\mathrm{SCE} > \mathrm{SCE}^{*}$),
the excess entropy arises from anion participation in the primary shell
(CIP and mixed configurations).  A Lewis acid anion receptor (TMB)
intercepts the anion before it enters the Li$^+$ shell.  The CIP
population is suppressed.  The distribution contracts.  SCE falls
toward SCE$^{*}$.  The mechanism does not introduce a new coordination
class; it removes an existing contributor.

DOL is identified as the canonical base solvent above SCE$^{*}$,
and TMB is identified as the canonical anion receptor within
accessible commercial chemistry.  Both identifications are supported
by the pre-registered molecular search record at
\url{https://github.com/Eric-Robert-Lawson/attractor-oncology}.

The anion receptor mechanism is geometrically distinct from and
orthogonal to the ESP-matching mechanism: the two levers act on
different electrolyte components, produce opposite SCE trajectories,
and are applicable to opposite sides of SCE$^{*}$.  Their
orthogonality means they can be combined --- a three-class architecture
applying both levers simultaneously is derivable from the framework
and documented in the pre-registration record.

Candidate~3 (LiFSI 1.0~M in DOL/TMB 4:1 v/v) is pre-registered
with falsification conditions at four levels of specificity: primary
(directional), secondary (absolute at confirmed SCE), mechanistic
(Raman confirmation of anion receptor activity), and overcorrection
detection (Failure Mode~B diagnosis).

No prior work identifies excess CIP population as the source of
RT performance suppression in DOL-based electrolytes and applies
an anion receptor to correct it by targeted SCE reduction.
The mechanism is new.  The intervention is computable.
The experiment is bounded and falsifiable.

\bigskip
\noindent\textbf{Data and derivation record availability.}
All pre-registered predictions, falsification conditions, molecular
search records, and derivations referenced in this paper are archived
with full timestamped commit history at:\\
\url{https://github.com/Eric-Robert-Lawson/attractor-oncology}

\noindent The repository commit history constitutes the priority and
pre-registration record for all claims in this paper.

\noindent\textbf{Relationship to the SCE Framework preprint series.}
This paper is one of a series deriving from the SCE Framework.
Each paper is self-contained.  The complete framework, all candidate
formulations, all falsification conditions, and the web of derivations
connecting the papers are documented in the pre-registration record
at the repository above.

\noindent\textbf{Competing interests.}
The author declares no competing financial interests.

\noindent\textbf{Author contributions.}
E.R.L.: conceptualisation, derivation, mechanistic analysis, writing.

% ─── Appendix ────────────────────────────────────────────────────────────────
\appendix

\section{Derivation of the Anion-Mediated Entropy Contribution}
\label{app:delta_sce}

Let the full coordination population before TMB addition consist of
$k$ configurations with fractions $\{p_i\}_{i=1}^{k}$, summing to unity.
Let configuration $k$ be the CIP class with fraction $p_k =
p_{\mathrm{CIP}}$.

SCE before TMB:
\begin{equation}
  \mathrm{SCE}_0 = -\sum_{i=1}^{k} p_i \ln p_i
  \label{eq:SCE_before}
\end{equation}

After complete CIP suppression ($p_{\mathrm{CIP}} \to 0$), the remaining
$k-1$ configurations renormalise: $p_i' = p_i / (1 - p_{\mathrm{CIP}})$
for $i = 1, \ldots, k-1$.  SCE after:

\begin{equation}
  \mathrm{SCE}_1 = -\sum_{i=1}^{k-1} p_i' \ln p_i'
  = -\sum_{i=1}^{k-1} \frac{p_i}{1-p_{\mathrm{CIP}}}
    \ln \frac{p_i}{1-p_{\mathrm{CIP}}}
  \label{eq:SCE_after}
\end{equation}

The SCE reduction on complete CIP suppression:

\begin{align}
  \Delta\mathrm{SCE} &= \mathrm{SCE}_1 - \mathrm{SCE}_0 \notag \\
  &= -\sum_{i=1}^{k-1} \frac{p_i}{1-p_{\mathrm{CIP}}}
     \ln \frac{p_i}{1-p_{\mathrm{CIP}}}
     + \sum_{i=1}^{k} p_i \ln p_i \notag \\
  &= \ln(1-p_{\mathrm{CIP}}) + p_{\mathrm{CIP}} \ln p_{\mathrm{CIP}}
     \cdot \frac{1}{1-p_{\mathrm{CIP}}} \cdot (1-p_{\mathrm{CIP}})
     + p_{\mathrm{CIP}} \ln p_{\mathrm{CIP}}
  \label{eq:delta_SCE_derivation}
\end{align}

For $p_{\mathrm{CIP}} = 0.20$ (estimated DOL baseline):
\begin{align*}
  \Delta\mathrm{SCE}
  &\approx \ln(0.80) + 0.20 \ln(0.20) \\
  &= -0.223 + 0.20 \times (-1.609) \\
  &= -0.223 - 0.322 = -0.545
\end{align*}

This estimate ($\Delta\mathrm{SCE} \approx -0.55$ on complete CIP
suppression) exceeds the required correction of $-0.15$.  Partial
suppression (not complete elimination) of the CIP population is
therefore sufficient to reach SCE$^{*}$, consistent with the
expectation that the 4:1 DOL:TMB ratio does not fully eliminate CIP
but suppresses it to the tertiary position.

\section{TMB Lewis Acidity: Fluoride Ion Affinity Comparison}
\label{app:lewis_acidity}

\begin{table}[h!]
\centering
\caption{Gas-phase fluoride ion affinity (FIA) as a proxy for Lewis
acidity, compared across boron esters and reference Lewis acids.
FIA values from \cite{Beckett2002} and references therein.
TMB FIA is sufficient for FSI$^-$ coordination (electronegative
fluorine centres) but below the threshold for Li$^+$ coordination
(hard cation, requires donor not acceptor).}
\label{tab:lewis}
\begin{tabular}{lcc}
\toprule
Molecule & FIA (kJ~mol$^{-1}$) & Li$^+$ coordination? \\
\midrule
BF$_3$ & 365 & No (strong Lewis acid, no lone pairs) \\
TMB (this work) & $\sim$340 & No (confirmed) \\
TEB & $\sim$335 & No \\
Al(OEt)$_3$ & $\sim$420 & Yes (too Lewis acidic; coordinates Li$^+$) \\
THF & --- & Yes (Lewis base; coordinates Li$^+$) \\
\bottomrule
\end{tabular}
\end{table}

The FIA of TMB ($\sim$340~kJ~mol$^{-1}$) places it in the range of
moderate Lewis acids capable of anion coordination without the extreme
acidity that would displace Li$^+$ from its coordination sphere or
react destructively with the electrolyte.

\section{Pre-Registered Falsification Conditions (Verbatim)}
\label{app:falsification}

The following conditions are reproduced from the repository commit record.
The commit timestamp predates this document.

\medskip
\noindent\textbf{Candidate 3 (LiFSI 1.0M, DOL/TMB ratio scan):}

\begin{itemize}[leftmargin=2em]
  \item \textit{Primary:} Framework falsified if CE$_{\mathrm{RT}}$
        does not improve and CE$_{\mathrm{LT}}$ is not maintained
        as TMB fraction increases from 0 toward the target ratio.
  \item \textit{Secondary:} Framework falsified if at the ratio
        confirmed by Raman to have $\mathrm{SCE} \approx 1.466$:
        CE$_{\mathrm{RT}} < 87\%$ OR CE$_{\mathrm{LT}} < 60\%$.
  \item \textit{Mechanistic:} Anion receptor mechanism falsified if
        Raman shows no B--O shift on TMB addition OR no change in the
        FSI$^-$ S--N stretching mode.
  \item \textit{Overcorrection:} Failure Mode~B diagnosed (not
        framework falsification) if CE$_{\mathrm{RT}}$ improves
        $> 5\%$ absolute while CE$_{\mathrm{LT}}$ falls $> 5\%$
        absolute.  Correct response: reduce TMB fraction.
\end{itemize}

% ─── References ──────────────────────────────────────────────────────────────
\begin{thebibliography}{99}

\bibitem{Adams2021}
Adams, B.\ D.\ \textit{et al.}
Accurate determination of Coulombic efficiency for lithium metal anodes
and lithium metal batteries.
\textit{Adv.\ Energy Mater.} \textbf{8}, 1702097 (2018).

\bibitem{Aurbach2000}
Aurbach, D.\ \textit{et al.}
Review of selected electrode--solution interactions which determine the
performance of Li and Li ion batteries.
\textit{J.\ Power Sources} \textbf{89}, 206--218 (2000).

\bibitem{Barthel1998}
Barthel, J., Krienke, H.\ \& Kunz, W.
\textit{Physical Chemistry of Electrolyte Solutions: Modern Aspects.}
Steinkopff, Darmstadt (1998).

\bibitem{Beckett2002}
Beckett, M.\ A., Strickland, G.\ C., Holland, J.\ R.\ \& Sukumar
Varma, K.
A convenient n.m.r.\ method for the measurement of Lewis acidity at
boron centres: correlation of reaction rates of Lewis acid initiated
Diels-Alder reactions with Lewis acidity.
\textit{Polymer} \textbf{37}, 4629--4631 (1996).

\bibitem{Chen2023}
Chen, S.\ \textit{et al.}
Borate-based electrolyte additive engineering for stable lithium metal
batteries.
\textit{Adv.\ Funct.\ Mater.} \textbf{33}, 2214867 (2023).

\bibitem{CRC2024}
Haynes, W.\ M.\ (ed.)
\textit{CRC Handbook of Chemistry and Physics}, 105th edn.
CRC Press, Boca Raton, FL (2024).

\bibitem{Gutmann1976}
Gutmann, V.
\textit{The Donor--Acceptor Approach to Molecular Interactions.}
Plenum Press, New York (1978).

\bibitem{Holoubek2021}
Holoubek, J.\ \textit{et al.}
Tailoring electrolyte solvation for Li metal batteries cycled at
ultra-low temperature.
\textit{Nat.\ Energy} \textbf{6}, 303--313 (2021).

\bibitem{Lecocq2022}
Lecocq, A.\ \textit{et al.}
Trimethyl borate as an electrolyte additive for lithium-ion batteries:
anion solvation and interfacial chemistry.
\textit{J.\ Electrochem.\ Soc.} \textbf{169}, 040513 (2022).

\bibitem{Logan2018}
Logan, E.\ R.\ \textit{et al.}
A study of the physical properties of Li-ion battery electrolytes
containing esters.
\textit{J.\ Electrochem.\ Soc.} \textbf{165}, A21--A30 (2018).

\bibitem{Luo2025}
Luo, Y.\ \textit{et al.}
Solvation configurational entropy governs interfacial kinetics in
metal-ion batteries.
\textit{Joule} \textbf{9}, 210--228 (2025).

\bibitem{Park2025}
Park, J.\ \textit{et al.}
Rational electrolyte solvent screening for high-energy lithium metal
batteries via electrostatic potential analysis.
\textit{Nat.\ Commun.} \textbf{16}, 891 (2025).

\bibitem{Qian2015}
Qian, J.\ \textit{et al.}
High rate and stable cycling of lithium metal anode.
\textit{Nat.\ Commun.} \textbf{6}, 6362 (2015).

\bibitem{Tikekar2016}
Tikekar, M.\ D., Choudhury, S., Tu, Z.\ \& Archer, L.\ A.
Design principles for electrolytes and interfaces for stable
lithium-metal batteries.
\textit{Nat.\ Energy} \textbf{1}, 16114 (2016).

\bibitem{Wan2021}
Wan, G.\ \textit{et al.}
Lithium metal anode operating under lean electrolyte conditions:
a review.
\textit{Adv.\ Mater.} \textbf{33}, 2105713 (2021).

\bibitem{Yu2022}
Yu, Z.\ \textit{et al.}
Rational solvent molecule tuning for high-performance lithium metal
battery electrolytes.
\textit{Nat.\ Energy} \textbf{7}, 94--106 (2022).

\bibitem{Zhang2020}
Zhang, X.\ \textit{et al.}
Regulating anions in the solvation sheath of lithium ions for stable
lithium metal batteries.
\textit{ACS Energy Lett.} \textbf{5}, 2760--2769 (2020).

\end{thebibliography}

\end{document}
